\section{Introduction}
\label{sec:intro}

\begin{intuition}
\textbf{In one paragraph.} When an embedding-based memory system
replaces many items with a few representatives, what determines how
much \emph{identity} can survive that compression? This paper proves
that for any consolidator acting on a unit-norm embedding cluster, the
answer is a single spectral ratio: the within-cluster spread $\dbar$
and the retrieval cap $\thetap$, raised to half the cluster's
effective dimension $\deff$. Two concrete things follow. One, there is
a phase boundary at $\dbar = \thetap$: below it every strategy
works; above it strategies diverge by tens of identity points in an
order the law predicts. Two, the law's main practical consequence on
real text is not that an adaptive method wins. It is that simple
\emph{centroid} averaging is already near-optimal --- because real
text clusters tend to sit in exactly the tight regime where the
theorem says averaging is safe. The bound is the same spectral
quantity that governs retrieval under noise; a bound on \emph{losing}
information is also a bound on \emph{compressing} it.
\end{intuition}

\subsection{The question}
\label{sec:intro:question}

Consider an embedding-based vector store that has indexed a million
text items, each mapped to a unit-length vector by a sentence encoder.
A query arrives and the store must return the items most similar to
the query under cosine similarity. At production scale one cannot
keep all million vectors at full precision: related items are
\emph{consolidated} into a smaller set of representatives, and only
the representatives are ranked at query time.

Consolidation takes many forms. A cluster's arithmetic
\emph{centroid} can be used as its representative; its \emph{medoid}
(the real item closest to the rest) can be used instead; the cluster
can be \emph{pruned} to its densest half; or an adaptive router can
pick among these per-cluster from local geometric features. Each
choice has its defenders. None of them by itself tells a designer
how much identity has actually been lost.

The scientific question is sharper than which choice is best:
\emph{what geometric quantity of a cluster determines how much
identity can survive consolidation, and does that quantity predict
which operator is appropriate?} If that quantity exists, it should
impose a floor on identity loss that no choice of operator can
escape; it should separate a safe regime from an unsafe one; and it
should explain, rather than merely record, the performance of each
operator on real text.

\subsection{The answer}
\label{sec:intro:answer}

We prove that such a quantity exists. For any unit-norm cluster
$\clusterset \subset S^{d-1}$ with mean within-cluster cosine distance
$\dbar$ and local effective dimension $\deff$ (participation ratio of
its covariance spectrum), and for any consolidator mapping the
cluster to $m$ representatives, whenever the retrieval cap $\thetap
:= 1-\theta$ is smaller than $\dbar$,
\[
\eps_{\mathrm{id}} \;\geq\; 1 - c_1\,m\,(\thetap/\dbar)^{\deff/2},
\]
where $c_1$ is a universal constant (Theorem~\ref{thm:duality}). We
call this the \emph{Consolidation--Interference Duality}: the
identical spectral quantity that sets an information-theoretic floor
on retrieval under noise~\cite{priceOfMeaning2026} reappears as a
floor on consolidation under compression. The bound is a direct
union-bound application of a cap-volume lemma; what is new is that
it binds compression, not just retrieval, and that it binds it with
the local effective dimension of the cluster rather than the ambient
dimension.

The bound separates two regimes. In the \emph{tight regime}
$\dbar < \thetap$ the bound vanishes and every operator has room to
preserve identity; experimentally, every strategy achieves mean
cap-coverage error $\leq 0.5\%$ on tight synthetic cells. In the
\emph{spread regime} $\dbar \geq \thetap$ the bound bites and
operators diverge; experimentally, mean cap-coverage errors diverge
to $30$--$74\%$ with an order uniquely predicted by $\deff$. This
tight/spread split is not a soft trend. It is visible as a phase
boundary in $(\deff, \theta)$ space and organises the ordering of
real-corpus identity accuracy across six sentence encoders.

\subsection{The surprising consequence}
\label{sec:intro:consequence}

If the law says the tight regime is where identity is preserved, a
natural next question is: where do real-text clusters sit? The
answer we find, across six corpora (synthetic DRM, MS~MARCO, Natural
Questions, HotpotQA, Wikipedia sections, arXiv titles) and six
contrastive sentence encoders, is that they sit mostly in the tight
regime --- and in that regime the geometry is lenient enough that
simple centroid averaging is already near-optimal. Three
experimental findings make this claim sharp.

\begin{enumerate}[leftmargin=1.8em, itemsep=2pt]
\item An adaptive per-cluster router we built (Geometry-Aware
Consolidation, GAC) Pareto-dominates the medoid baseline on every
synthetic cell, but a \emph{fixed-centroid} ablation beats the full
router by $1$--$6$ identity points on every one of the five
real-text corpora (\S\ref{sec:gac:ablation}). The residual-direction
budget that dominates the router's engineering complexity contributes
nothing ($\Delta \leq 0.002$). A per-cluster \emph{oracle} barely
improves over fixed centroid on real text.
\item At matched bytes-per-vector against five learned-quantization
baselines (product quantization~\cite{jegouPQ2011}, optimized product
quantization~\cite{ge2013optimized}, LSH~\cite{charikarLSH2002},
PCA+int8, HNSW-prune~\cite{malkovHNSW2018}), consolidation and
learned quantization populate different Pareto corners of the
identity-vs-bytes frontier; centroid dominates on the five
low-to-moderate-$\deff$ corpora and quantization takes over only on
high-$\deff$ arXiv titles (\S\ref{sec:e3}).
\item Closing the loop to a downstream reader, a
Llama-3.1-70B-Instruct retrieval-augmented pipeline on three QA
benchmarks shows a regime-dependent three-way split: centroid
\emph{hurts} Natural Questions EM by $4.2$ points, is neutral on
HotpotQA, and \emph{wins} by $8.4$ EM points on PopQA --- matching
the cap-coverage prediction of which operator the law prefers on
each corpus (\S\ref{sec:e7}).
\end{enumerate}

The centroid result is not an engineering accident; it is the
law's main practical implication in the regime where real text
lives. The residual-direction and adaptive-routing machinery that
the engineering literature builds in anticipation of a harder
regime turns out, on the corpora we tested, to be solving a
problem that does not arise.

\subsection{GAC is a probe, not the main product}
\label{sec:intro:gac_as_probe}

A paper whose primary product is an adaptive router would want that
router to win. This paper uses GAC for a different purpose: as an
\emph{instrument} that asks whether, given the regime structure the
theorem predicts, adaptive per-cluster routing can outperform a
fixed simple operator. The instrument returns a clean negative
answer on real text. The negative is scientifically informative ---
it localises the value of adaptation to the low-$\dbar$ synthetic
regime, it isolates which components of the router matter
(none of the residual-direction machinery does), and it sets up the
downstream three-way finding. We therefore describe GAC in full
(\S\ref{sec:gac}), but the paper is organised so its role is to
test the law rather than to advertise a method.

\subsection{Scope of the law}
\label{sec:intro:scope}

The theorem is stated and proved for unit-norm embedding clusters
under cosine-threshold retrieval. Its experimental support is
text-centric: six contrastive sentence encoders on six text corpora,
at scales from $10^3$ to $10^6$ items. Under those assumptions the
bound is provably a lower bound on identity-retrieval error, and
empirically it organises the behaviour of every consolidation
operator we tested. What the paper does \emph{not} claim is a full
theory of cognitive memory or of arbitrary learned-state memory in
neural networks. Biological consolidation, multimodal embedding
memory, and replay in large-state continual-learning agents share
the same \emph{structural} compression problem and so are natural
settings in which to test whether the law extends; but the
evidence we report is for text embedding consolidation, and we keep
the scientific claim inside that envelope.

The theorem itself is also honest about where it is strongest. In
the tight regime (defined directly by the theorem's hypothesis
$\dbar < \thetap$) the calibrated constant sits at $c_1^{95}
\approx 0.05$ on $12{,}400$ cells --- within an order of magnitude
of the ideal Berry--Esseen constant that falls out of the
cap-volume lemma; in the spread regime, the isotropic cap-volume
approximation is loose and the calibrated $c_1$ grows unbounded as
$\theta \to 1$ at small $\deff$.
We state the theorem as the first mathematically explicit version of
the law and mark an \emph{anisotropic refinement} of the cap-volume
argument as the central open theoretical problem
(\S\ref{sec:limits}).

\subsection{Where this sits in a trilogy}
\label{sec:intro:trilogy}

This paper is the third of three on the geometry of embedding
memory, and can be read independently of the other two. The first,
\emph{The Geometry of Forgetting}~\cite{geometryOfForgetting2026},
documented empirically that production sentence encoders concentrate
$\geq 99\%$ of their variance in at most $\deff \approx 16$ effective
dimensions, across six encoder families and ten benchmarks. The
second, \emph{The Price of Meaning}~\cite{priceOfMeaning2026}, lifted
that observation to a theorem: any kernel-threshold memory with
finite $\deff$ must forget by interference under retrieval noise,
with a lower bound no storage budget can evade. The present paper
closes the trilogy by showing that the compression lever does not
provide an escape route: the same spectral quantity that governs
forgetting governs consolidation, by the same cap-volume argument
applied to the cluster-conditional measure. The three papers share
a mathematical core; what they do not share is the object under
compression (global store vs.\ cluster) or the failure mode under
study (retrieval noise vs.\ compression).

\subsection{Contributions}
\label{sec:intro:contributions}

In order of scientific centrality:

\begin{description}[leftmargin=1.8em, itemsep=3pt]
\item[\textbf{(C1)} The Duality Theorem.]
A universal-constant lower bound on identity-retrieval error after
consolidation, in the same spectral quantity
$(\thetap/\dbar)^{\deff/2}$ that governs retrieval under noise
(\S\ref{sec:duality}; proof in Appendix~\S\ref{sec:methods:proof}).

\item[\textbf{(C2)} Regime structure and empirical validation.]
The theorem induces a tight/spread phase boundary at $\dbar =
\thetap$. We validate the boundary on a $400$-cell synthetic grid
($16{,}000$ trials), and show the same regime structure organises
the ordering of six real-text corpora across six sentence encoders
(\S\ref{sec:e1_validation}, \S\ref{sec:template_collapse}).

\item[\textbf{(C3)} Centroid dominates on real text.]
On five real-text corpora, a fixed-centroid operator dominates a
stochastic adaptive router by $1$--$6$ identity points; the
residual-direction budget contributes nothing; a per-cluster oracle
barely improves. This is not an engineering accident --- it is the
law's prediction for corpora that sit in the tight regime
(\S\ref{sec:gac:ablation}).

\item[\textbf{(C4)} Geometry selects between consolidation and
learned quantization.] At matched bytes-per-vector, consolidation
and learned quantization occupy different Pareto corners of the
identity-vs-bytes frontier. The crossover is predicted by $\deff$:
quantization takes over on high-$\deff$ arXiv titles; consolidation
dominates on the other five corpora (\S\ref{sec:e3}).

\item[\textbf{(C5)} Downstream consequences are regime-dependent.]
A Llama-3.1-70B-Instruct retrieval-augmented pipeline on Natural
Questions, HotpotQA, and PopQA exhibits a regime-dependent
three-way split in end-to-end EM, matching the cap-coverage
prediction of which operator the law prefers per corpus
(\S\ref{sec:e7}).

\item[\textbf{(C6)} Limits of the current theory.]
The isotropic cap-volume bound is shape-correct and predictively
useful in the tight regime ($c_1^{95} \approx 0.05$, $95\%$
coverage; $c_1^{99} \approx 370$, $99\%$) and loose-to-vacuous in
the spread regime ($c_1 \to
\infty$ as $\theta \to 1$ at small $\deff$). An anisotropic refinement of the
cap-volume argument is the central open theoretical problem
(\S\ref{sec:limits}).

\item[\textbf{(C7)} GAC as an instrument.]
We introduce Geometry-Aware Consolidation and use it as a probe of
what per-cluster adaptation can buy under the law. GAC
Pareto-dominates the medoid baseline on every synthetic cell. Its
failure to beat fixed centroid on real text localises the value of
adaptation to the low-$\dbar$ synthetic regime and is itself one of
the paper's scientific findings (\S\ref{sec:gac}).
\end{description}

\subsection{How to read this paper}

We have split the exposition so that different readers can start at
different depths.

\begin{itemize}[leftmargin=1.8em, itemsep=2pt]
\item A reader who wants the \textbf{law only} can stop after
Section~\ref{sec:duality} and skim the regime-map figure in
Section~\ref{sec:e1_validation}.
\item A reader who wants the \textbf{centroid lesson on real text}
should read Sections~\ref{sec:template_collapse} and~\ref{sec:gac:ablation}.
\item A reader who wants the \textbf{downstream implications for RAG}
should jump to Section~\ref{sec:e7}.
\item A reader who wants the \textbf{proof} should read
Section~\ref{sec:duality} then Appendix~\S\ref{sec:methods:proof}.
\item A reader new to consolidation, effective rank, or RAG should
start with Section~\ref{sec:background}.
\end{itemize}

Each results section opens with a gray ``intuition'' box giving the
plain-English reading, followed by the full technical development.
The intent is that a reader who is technical but not steeped in
embedding geometry can read the intuition boxes end-to-end and leave
with the paper's full argument.
